%====================================
% KAPITEL 3: METHODISCHES VORGEHEN
%====================================
\section{Methodisches Vorgehen}

Das folgende Kapitel legt das methodische Vorgehen der Arbeit offen. Zur Beantwortung der Forschungsfragen wird eine systematische Literaturanalyse durchgef\"uhrt, deren Grundverst\"andnis und Ablauf zun\"achst begr\"undet werden (Kapitel~3.1). Anschlie\ss{}end werden die Suchstrategie mit den verwendeten Datenbanken und Schlagworten dokumentiert (Kapitel~3.2) sowie der Auswahlprozess nach dem PRISMA-Schema und die Auswertungslogik in Anlehnung an die qualitative Inhaltsanalyse nach Mayring dargelegt (Kapitel~3.3). Die vollst\"andige Offenlegung aller methodischen Entscheidungen dient der intersubjektiven Nachvollziehbarkeit und Replizierbarkeit der Untersuchung.

\subsection{Systematische Literaturanalyse nach Webster und Watson}

\textbf{Methodenwahl und Begr\"undung}

Die Wahl der systematischen Literaturanalyse als Forschungsmethode ergibt sich unmittelbar aus dem Erkenntnisziel der Arbeit. Die Forschungsfrage zielt nicht auf die Erhebung neuer Prim\"ardaten, sondern auf die Identifikation, Systematisierung und Bewertung bereits vorliegender empirischer Befunde zu Change-Management-Faktoren in ERP-Projekten des Mittelstands. F\"ur ein solches Erkenntnisziel ist die Literaturanalyse das Mittel der Wahl: Sie erm\"oglicht es, den fragmentierten Forschungsstand zweier bislang selten integrativ betrachteter Felder zusammenzuf\"uhren, Muster \"uber Einzelstudien hinweg sichtbar zu machen und Forschungsl\"ucken pr\"azise zu bestimmen.\footcite[Vgl.][S.~xv--xviii]{webster2002} Gegen\"uber einer eigenen empirischen Erhebung -- etwa Experteninterviews in einzelnen mittelst\"andischen Unternehmen -- bietet sie zudem den Vorteil einer breiteren empirischen Basis: Die Befunde vieler Studien mit unterschiedlichen Unternehmen, Branchen und L\"andern lassen sich verdichten, w\"ahrend eine Einzelerhebung im Rahmen einer Bachelorarbeit nur einen kleinen, kaum generalisierbaren Ausschnitt erfassen k\"onnte.

Entscheidend ist dabei die Abgrenzung zur unsystematischen, narrativen Literatur\"ubersicht: Eine systematische Literaturanalyse folgt einem vorab definierten, dokumentierten und damit replizierbaren Suchprozess mit expliziten Ein- und Ausschlusskriterien, wodurch die Gefahr einer selektiven, die eigenen Erwartungen best\"atigenden Quellenauswahl reduziert wird.\footcite[Vgl.][S.~xv--xviii]{webster2002} Sie erf\"ullt damit die G\"utekriterien wissenschaftlichen Arbeitens -- Objektivit\"at, Reliabilit\"at und Validit\"at -- in dem Ma\ss{}e, in dem der Prozess offengelegt und konsistent angewendet wird.\footcite[Vgl.][S.~882]{okoli2015}

\textbf{Das Vorgehen nach Webster und Watson}

Die Arbeit folgt dem methodischen Rahmen von Webster und Watson, der sich als Standard f\"ur Literaturanalysen in der Informationssystemforschung etabliert hat.\footcite[Vgl.][S.~xv--xviii]{webster2002} Zwei Grundprinzipien dieses Ansatzes pr\"agen das Vorgehen:

Erstens die mehrstufige Suchstrategie: Ausgehend von einer strukturierten Datenbanksuche in den f\"uhrenden Fachzeitschriften und Konferenzb\"anden wird der Quellenkorpus durch R\"uckw\"artssuche (Auswertung der Literaturverzeichnisse der identifizierten Beitr\"age) und Vorw\"artssuche (Identifikation j\"ungerer Beitr\"age, die die Kernquellen zitieren) erg\"anzt.\footcite[Vgl.][S.~xv--xviii]{webster2002} Dieses dreistufige Vorgehen stellt sicher, dass die Analyse nicht auf die Treffermenge einzelner Suchanfragen beschr\"ankt bleibt, sondern das Themenfeld m\"oglichst vollst\"andig erfasst -- ein Anspruch, der angesichts der verstreuten Publikationslandschaft an der Schnittstelle von ERP-, Change-Management- und Mittelstandsforschung besonders relevant ist.

Zweitens die konzeptzentrierte Auswertung: Webster und Watson pl\"adieren daf\"ur, die Literatur nicht autoren- oder studienweise, sondern entlang von Konzepten zu strukturieren.\footcite[Vgl.][S.~xv--xviii]{webster2002} F\"ur die vorliegende Arbeit liefert der in Kapitel~2.4 entwickelte Bezugsrahmen die konzeptionelle Struktur: Die identifizierten Beitr\"age werden entlang der Change-Management-Faktoren, der ADKAR-Elemente und der Erfolgsdimensionen des DeLone-\&-McLean-Modells ausgewertet und in einer Konzeptmatrix verdichtet, die die Grundlage der Ergebnisdarstellung in Kapitel~4 bildet.

\subsection{Suchstrategie, Datenbanken und Schlagworte}

Die Literatursuche erfolgt in neun wissenschaftlichen Datenbanken, die f\"ur die Wirtschaftsinformatik und angrenzende Disziplinen relevant sind: SpringerLink, ScienceDirect, IEEE Xplore, die ACM Digital Library, Wiley Online Library, Web of Science, EBSCOhost, die AIS Electronic Library (AISeL) sowie das f\"ur deutschsprachige Wirtschaftsliteratur einschl\"agige Portal WISO-net. Der Suchalgorithmus kombiniert Boolesche Operatoren, um die zentralen Konstrukte der Arbeit pr\"azise abzudecken. Die Suchstring-Formulierung lautet:

\begin{quote}
\verb|("ERP implementation" OR "Enterprise Resource Planning")|\\
\verb|AND ("change management")|\\
\verb|AND ("critical success factors" OR "success factors")|\\
\verb|AND ("SME" OR "Mittelstand")|
\end{quote}

F\"ur die deutschsprachige Suche wird eine \"aquivalente Variante verwendet:

\begin{quote}
\verb|("ERP-Einf"uhrung" OR "ERP-Implementierung")|\\
\verb|AND ("Change Management" OR "Ver"anderungsmanagement")|\\
\verb|AND ("Erfolgsfaktoren")|\\
\verb|AND ("Mittelstand" OR "KMU")|
\end{quote}

Die Qualit\"atssicherung der identifizierten Literatur erfolgt \"uber mehrere Stufen. Ber\"ucksichtigt werden ausschlie\ss{}lich peer-reviewte Beitr\"age aus wissenschaftlichen Zeitschriften und Konferenzb\"anden sowie etablierte Fachb\"ucher. Die Zeitschriftenqualit\"at wird anhand des H-Index und der Einstufung im VHB-JOURQUAL-Ranking bewertet. Um die wissenschaftliche Fundierung der Arbeit sicherzustellen, wird ein ausgewogenes Verh\"altnis zwischen Prim\"ar- und Sekund\"arquellen angestrebt: Der Anteil peer-reviewter Journalartikel und Konferenzbeitr\"age soll mindestens 60~\% des Quellenkorpus betragen.

\subsection{Auswahlprozess nach PRISMA und Auswertungslogik nach Mayring}

Der Auswahlprozess der identifizierten Literatur folgt dem PRISMA-Statement (Preferred Reporting Items for Systematic Reviews and Meta-Analyses).\footcite[Vgl.][n71]{page2021} Dieser mehrstufige Prozess umfasst die Phasen Identifikation, Screening, Eignungspr\"ufung und finalen Einschluss der Studien. Um die Transparenz des Auswahlprozesses zu erh\"ohen, wird ein PRISMA-Flowchart erstellt, das die Anzahl der identifizierten, gescreenten und final eingeschlossenen Publikationen dokumentiert.

Gegenstand der Untersuchung ist die aktuelle Forschungsliteratur zu Change-Management-Faktoren bei ERP-Implementierungen, mit einem Schwerpunkt auf Publikationen ab dem Jahr 2015. Diese zeitliche Eingrenzung stellt sicher, dass der technologische und organisatorische Wandel der vergangenen Jahre -- insbesondere die zunehmende Cloud-Nutzung und agile Implementierungsmethoden -- in den analysierten Studien abgebildet wird. Die angestrebte Datenbasis umfasst voraussichtlich etwa 20 bis 30 peer-reviewte Quellen. Die geografische Eingrenzung auf den deutsch- und englischsprachigen Raum gew\"ahrleistet sowohl die Vergleichbarkeit der wirtschaftlichen Rahmenbedingungen als auch die Breite des internationalen Forschungsstands.

Die Auswertung der identifizierten Faktoren erfolgt mittels einer inhaltlich strukturierenden Inhaltsanalyse nach Mayring.\footcite[Vgl.][S.~50~ff.]{mayring2015} Dabei werden die in der Literatur genannten Change-Management-Faktoren zun\"achst extrahiert, anschlie\ss{}end induktiv und deduktiv zu Kategorien verdichtet und schlie\ss{}lich hinsichtlich ihrer Relevanz f\"ur den Mittelstand bewertet. Als theoretischer Rahmen f\"ur die Kategorienbildung dienen die zwei etablierten Modelle der Arbeit: das IS-Erfolgsmodell nach DeLone und McLean\footcite[Vgl.][S.~13--16]{delone2003} sowie das ADKAR-Modell.\footcite[Vgl.][S.~126--127]{galli2018}
